\centerline{\bf \Large  10 класс}

\problem{\textbf{1}} Пюрвя загадал число, которое имеет $2025$ различных делителей, отличных от $1$ и самого числа. Сколько простых делителей может быть у числа Пюрви?

\textbf{Решение:} Заметим, что у числа всего 2027 делителей, считая 1 и само число. Вспомним, как считается количество делителей числа:
если степени вхождения простых чисел в его разложение равны $a_1, a_2, ..., a_n$, то количество делителей равно
$(a_1+1)(a_2+1)...(a_n+1)=2027$.

Но заметим, что 2027 простое число, а значит, в произведении участвует лишь один множитель, то есть число является 2026-ой степенью некоторого простого числа, откуда делаем вывод, что число может иметь только 1 простой делитель.


\textbf{Критерии.} \\
\textit{7 баллов} $-$ Полностью верное решение\\
\textit{0-6 баллов} $-$ Полностью верное решение, но учитывается само число и единица (делителей 2025)\\
\textit{0-6 баллов} $-$ Все остальные продвижения\\
\textit{0 баллов} $-$ Любой ответ без объяснений\\

\problem{\textbf{2}} Решите задачу:

а) Арслан бежит из Элисты в Оргакин. Он тратит 25\% времени на путь в гору, 60\% — по равнине, а остальное время — с горы. Время его движения из Элисты в Оргакин и по той же дороге из Оргакина в Элисту одинаково, а его скорости в гору, с горы и по равнине постоянны, но различны. Во сколько раз быстрее Арслан бежит с горы, чем в гору?

б) Составьте и решите обратную задачу по схеме: ?; 60; $\dfrac{5}{3}$.

\textbf{Решение А:}

\textbf{ВАЖНО! Доказать равенство длины подъёма и длины спуска, если оно используется в решении.}

Пусть скорость Арслана с горы в $x \neq 1$ раз больше скорости в гору.

Подъём на первом участке (слева направо) занимал $0,25$ от времени пути, тогда спуск на первом участке (справа налево) занимает $0,25 / x$ - так как скорость во время спуска в $x$ раз больше, следовательно, время в $x$ раз меньше.

Равнина - второй участок.

Спуск на третьем участке (слева направо) занимал $0,15$ от времени пути, тогда подъём на третьем участке занимает $0,15 x$ - так как скорость во время подъёма в $x$ раз меньше, следовательно, время в $x$ раз больше.

$$
\begin{aligned}
0,25+0,6+0,15= & 0,25 / x+0,6+0,15 x \quad \\
& \Rightarrow \quad 0,15 x^2-(0,25+0,15) x+0,25=0 \\
& \Rightarrow(x-1)(x-5 / 3)=0 \quad \Rightarrow \quad x=5 / 3 .
\end{aligned}
$$


\textbf{Примерное условие Б:}  Арслан бежит из Элисты в Оргакин. Он тратит некоторое количество времени на путь в гору, 60\% — по равнине, а остальное время — с горы. Время его движения из Элисты в Оргакин и по той же дороге из Оргакина в Элисту одинаково, а его скорости в гору, с горы и по равнине постоянны, но различны. Скорость в гору в $\dfrac{5}{3}$ раз больше скорости с горы. Какой процент от общего времени Арслан тратит на путь в гору из Элисты в Оргакин?

\textbf{Решение Б:} Пусть некоторый процент $100A$ он тратил на путь c горы, тогда если скорость Арслана с горы в $\dfrac{5}{3}$ раз больше скорости в гору, то из условия имеем:

$$
\begin{aligned}
0,4-A+A= & (0,4-A )/ \dfrac{5}{3}+A \dfrac{5}{3} \quad \Rightarrow \quad A=0,15 .
\end{aligned}
$$

Тогда на путь в гору тратит 40-15=25\%.

\textbf{Критерии.} \\
\textit{0-4 баллов} $-$ Решение пункта А\\
\textit{0-1 балла} $-$ Составление условия пункта Б \\
\textit{0-2 баллов} $-$ Решение пункта Б\\
\textit{0 баллов} $-$ Любой ответ без объяснений\\

\problem{\textbf{3}} Два приведённых квадратных трёхчлена $f(x)$ и $g(x)$ таковы, что каждый из них имеет по два корня, и выполняются равенства $f(2024)=g(2025)$ и $g(2024)=f(2025)$. Найдите сумму всех четырёх корней этих трёхчленов.

\textbf{Решение:} Пусть $f(x)=x^2+a x+b, g(x)=x^2+c x+d$. Тогда условия задачи запишутся в виде:

$$
2024^2+2024a+b=2025^2+2025 c+d \quad \text { и } \quad 2025^2+2025 a+b=2024^2+2024c+d .
$$


Вычитая из первого равенства второе, получаем $2024^2-2025^2-a=2025^2-2024^2+c$, то есть $a+c=-8098$. Но по теореме Виета $-a$ это сумма корней первого трёхчлена, а $-c$ сумма корней второго трёхчлена, откуда и следует, что сумма корней равна 8098.

\textbf{Критерии.} \\
\textit{7 баллов} $-$ Полностью верное решение\\
\textit{0-6 баллов} $-$ Все остальные продвижения\\
\textit{0 баллов} $-$ Любой ответ без объяснений\\

\problem{\textbf{4}} 
В остроугольном треугольнике $A B C$ провели высоты $B B_1, C C_1$. Прямая $B_1 C_1$ пересекает описанную окружность треугольника $A B C$ в точках $X$ и $Y$. Пуст $I_b$ - центр вписанной окружности $BXY$, а $I_c-C X Y$. Докажите, что треугольник $A I_b I_c$ - равнобедренный.

\textbf{Решение:} Докажем, что $A$ - середина дуги $X Y$, отсюда будет следовать утверждение задачи, ведь действительно: $I_b$ будет лежать на $B A$, $I_c$ на $C A$, применяя лемму о трезубце, $A I_b=A X=$ $A I_c$. Четырёхугольник $B C_1 B_1 C-$ вписанный на диаметре $B C$, откуда углы $A B C$ и $C B_1 Y$ равны (считаем, что $Y$ лежит на дуге АС - см. чертёж). Но если перевести это равенство в дуги, то получим, что сумма дуг $A Y$ и $Y C$ равна сумме дуг $A X$ и $Y C$, откуда и следует, что $A$ - середина дуги $X Y$.

\textbf{Критерии.} \\
\textit{7 баллов} $-$ Полностью верное решение\\
\textit{0-6 баллов} $-$ Все остальные продвижения\\
\textit{0 баллов} $-$ Любой ответ без объяснений\\

\problem{\textbf{5}} 
В Калмыкии есть 33 населённых пункта. Из каждого пункта выходит по 16 сухопутных и 16 речных транспортных путей, соединяющих его с остальными 32 пунктами (пути двусторонние). Назовём тройку населённых пунктов <<сухопутной>>, если среди них есть пункт, который соединён с двумя другими сухопутными путями, а путь между этими двумя пунктами — речной. Аналогично определим <<речную>> тройку. Докажите, что либо сухопутных, либо речных троек хотя бы 2025.

\textbf{Решение:} Определим \textit{вишню} как множество из двух различных рёбер графа $K_{33}$, имеющих общую вершину.

Замечаем, что монохроматический треугольник всегда содержит три монохроматические вишни, а полихроматический треугольник всегда содержит одну монохроматическую и две полихроматические вишни.

Поэтому рассматриваем величину $2 M-P$, где $M-$ количество монохроматических вишен, а $P$ - количество полихроматических вишен. Заметим, что каждая вишня является частью единственного треугольника, и распределим эту величину по всем различным треугольникам в $K_{43}$. По построению вклад полихроматического треугольника будет равен нулю, тогда как монохроматический треугольник даст вклад 6. Таким образом, получаем:

\centerline{$2 M-P=6 \cdot\{$ количество монохроматических треугольников $\}$.}


Рассмотрим произвольную вершину $v$. Пусть $M_v$-количество монохроматических вишен с центральной вершиной $v$, а $P_v$ - количество таких полихроматических вишен. Тогда справедливо:


$$
M_v=\frac{16 \cdot 15}{2}+\frac{16 \cdot 15}{2}=240 \quad \text { и } \quad P_v=16 \cdot 16=256 .
$$


Другими словами, для любой вершины $v$ выполняется:

$$
2 M_v-P_v=224.
$$


Суммируя все эти вклады, получаем:

$$
2 M-P= 33\cdot 402.
$$


Мы заключаем, что всего существует $\frac{33\cdot402}{6}=1232$ монохроматических треугольника. Поскольку все треугольников $C^3_{33} = 5456$, тогда полихроматических треугольников ровно $5456 - 1232 = 4224$. Из них хотя бы половина либо сухопутые, либо речные, то есть $\dfrac{4224}{2} = 2112 > 2025$.

\textbf{Критерии.} \\
\textit{7 баллов} $-$ Полностью верное решение\\
\textit{0-6 баллов} $-$ Все остальные продвижения\\
\textit{0 баллов} $-$ Любой ответ без объяснений\\